\section{Adaptive oracle theorem}
\label{sec:adaptive-oracle}

This section joins the finite-sample current-risk theory with a single adaptive
algorithm.  The result is an average current-risk oracle inequality.  It does not
claim pointwise recovery immediately after an unrestricted jump, which is
impossible under the model of Section~6.1.

\subsection{A finite library of predictable neural experts}

Let \(\mathcal I_T\) be a finite collection of deterministic hyperparameter
indices \(i=(m,H,g)\).  At prediction time \(t\), expert \(i\) fits the bounded
neural estimator of Section~1 using only \(s<t\) and deterministic temporal
weights
\[
 w^{H,g}_{t,s}
 \propto
 \left[1-\left(\frac{t-s}{H+1}\right)^g\right]_+,
 \qquad 0<g\le 1,
\]
with the obvious truncation by the available history.  Other deterministic
regular kernels may be added to the library.  Denote the resulting
past-measurable forecast by \(\widehat f_{t,i}\), and assume
\(\|\widehat f_{t,i}\|_\infty\le B\).

Let \(\mathcal F_{t-1}\) be the sigma-field generated by observations strictly
before \(t\).  Every \(\widehat f_{t,i}\) is \(\mathcal F_{t-1}\)-measurable.
Thus the experts may themselves be re-estimated from the past; the prediction
with expert advice argument does not require static experts.

\subsection{A hierarchical fixed-share master}

Put
\[
 z_T=\log(8T/\delta),\qquad
 u_T=\sqrt{2\sigma^2 z_T}+2bz_T,\qquad
 D_T=2B+u_T.
\]
Under Assumption~1.1,
\[
 \Pr\!\left(\max_{t\le T}|\varepsilon_t|>u_T\right)\le \delta/4.
\]
On the complementary event, \(|Y_t-\widehat f_{t,i}(X_t)|\le D_T\) for every
expert.  Define the normalized square loss
\[
 \ell_t(a)=\frac{(Y_t-a)^2}{D_T^2}.
\]
For \(a,Y_t\) in the preceding ranges, \(a\mapsto
\exp\{-\eta\ell_t(a)\}\) is concave whenever \(\eta\le 1/2\).  Indeed,
the second derivative is nonpositive because
\(2\eta(Y_t-a)^2/D_T^2\le 1\).

Fix a prior \(\pi\) on \(\mathcal I_T\).  For a share parameter
\(q\in[0,1)\), define a Markov prior on expert paths by
\[
 \Pi_q(i_{1:T})
 =
 \pi_{i_1}
 \prod_{t=2}^T
 \left\{
 (1-q)\mathbf 1(i_t=i_{t-1})+q\pi_{i_t}
 \right\}.
\]
Let \(\mathcal Q_T\) be a finite grid of share parameters with prior \(\rho\).
The master is the hierarchical exponential aggregate with learning rate
\(\eta=1/2\): internally it runs one fixed-share recursion for every
\(q\in\mathcal Q_T\), and externally exponentially aggregates their issued
forecasts.  Equivalently, it aggregates expert paths under the hierarchical
prior
\[
 \Pi(i_{1:T})=\sum_{q\in\mathcal Q_T}\rho_q\Pi_q(i_{1:T}).
\]
Its forecast is denoted \(\widehat f_t^{\rm A}\).

\begin{lemma}[Pathwise mixable tracking]
\label{lem:mixable-path}
On the event \(\max_t|\varepsilon_t|\le u_T\), for every path
\(i_{1:T}\),
\[
 \sum_{t=1}^T
 \left[
 (Y_t-\widehat f_t^{\rm A}(X_t))^2
 -(Y_t-\widehat f_{t,i_t}(X_t))^2
 \right]
 \le
 2D_T^2\,\mathcal C_T(i_{1:T}),
 \tag{A.1}
\]
where
\[
 \mathcal C_T(i_{1:T})
 :=-\log \Pi(i_{1:T})
 \le
 \min_{q\in\mathcal Q_T}
 \left\{
 -\log\rho_q-\log\Pi_q(i_{1:T})
 \right\}.
\]
\end{lemma}

\paragraph{Proof.}
At each inner or outer aggregation step, exp-concavity gives
\[
 \ell_t\!\left(\sum_j p_j a_j\right)
 \le
 -2\log\sum_j p_j e^{-\ell_t(a_j)/2}.
\]
Telescoping the exponential weights recursion over the Markov path prior yields
\[
 \sum_t\ell_t(\widehat f_t^{\rm A})
 \le
 \sum_t\ell_t(\widehat f_{t,i_t})
 +2\{-\log\Pi(i_{1:T})\}.
\]
Multiplication by \(D_T^2\) proves the result. \(\square\)

For a uniform expert prior and a path with \(S\) switches, if the share grid
contains some \(q\in[S/(2(T-1)),\,2S/(T-1)]\subset(0,1/2]\), then
\[
 \mathcal C_T(i_{1:T})
 \lesssim
 \log|\mathcal I_T|
 +S\log\!\left(\frac{e|\mathcal I_T|T}{S}\right)
 +\log|\mathcal Q_T|.
 \tag{A.2}
\]
A dyadic grid
\[
 \mathcal Q_T=\{0,2^{-1},2^{-2},\ldots,2^{-\lceil\log_2T\rceil}\}
\]
has cardinality \(O(\log T)\) and supplies this approximation whenever
\(1\le S\le(T-1)/2\).  The case \(S=0\) is covered by \(q=0\).

\subsection{Fast conversion from realised loss to current integrated risk}

For a predictable comparator path \(i_{1:T}\), put
\[
 A_T=\sum_{t=1}^T\mathcal E_t(\widehat f_t^{\rm A}),
 \qquad
 B_T(i_{1:T})
 =\sum_{t=1}^T\mathcal E_t(\widehat f_{t,i_t}).
\]
The next lemma is the statistical bridge.

\begin{lemma}[Localized sequential Bernstein bridge]
\label{lem:bridge}
Under Assumption~1.1, there exist constants \(c_0,c_1<\infty\), depending
only on \(B,\sigma,b\), such that for every predictable path and every \(x>0\),
with probability at least \(1-e^{-x}\),
\[
 A_T-B_T
 \le
 \sum_{t=1}^T
 \left[
 (Y_t-\widehat f_t^{\rm A}(X_t))^2
 -(Y_t-\widehat f_{t,i_t}(X_t))^2
 \right]
 +c_0\sqrt{x(A_T+B_T)}+c_1x.
 \tag{A.3}
\]
\end{lemma}

\paragraph{Proof.}
Set \(d_t=\widehat f_t^{\rm A}-\widehat f_{t,i_t}\).  Conditional on
\(\mathcal F_{t-1}\), the centered difference between realised and population
losses is a martingale difference.  Expanding
\(Y_t=f_t^\star(X_t)+\varepsilon_t\) gives a centered bounded design term and
the noise term \(-2\varepsilon_td_t(X_t)\).  Since all forecasts and targets
are bounded by \(B\),
\[
 {\rm Var}(Z_t\mid\mathcal F_{t-1})
 \le
 (16B^2+4\sigma^2)\|d_t\|_{L^2(P_t)}^2.
\]
The conditional Bernstein moment assumption gives a uniform scale constant
depending only on \(B,b\).  A martingale Bernstein inequality therefore yields
\[
 -\sum_t Z_t
 \le
 c_0\sqrt{x\sum_t\|d_t\|_{L^2(P_t)}^2}+c_1x.
\]
Finally
\[
 \|d_t\|_{L^2(P_t)}^2
 \le
 2\mathcal E_t(\widehat f_t^{\rm A})
 +2\mathcal E_t(\widehat f_{t,i_t}),
\]
which proves (A.3) after changing \(c_0\). \(\square\)

The localization in (A.3) is essential.  A nonlocalized
\(O(\sqrt{T})\) conversion would destroy the faster nonparametric rates in
sufficiently smooth cases.

\begin{theorem}[Adaptive current-risk oracle inequality]
\label{thm:adaptive-oracle}
Suppose Assumptions~1.1--1.4 hold and every expert in
\(\mathcal I_T\) is the bounded estimator of Theorem~3.2, trained only on the
past.  Let the hierarchical fixed-share master above use a finite path prior
\(\Pi\).  Then, for every \(\zeta\in(0,1]\), with probability at least
\(1-\delta\),
\[
 \frac1T\sum_{t=1}^T
 \mathcal E_t(\widehat f_t^{\rm A})
 \le
 \inf_{i_{1:T}}
 \left\{
 (1+\zeta)\frac1T\sum_{t=1}^T
 \mathcal B_{t,i_t}
 +
 \frac{C_\zeta\,\Xi_T^2}{T}
 \left[
 \mathcal C_T(i_{1:T})+\log\frac1\delta
 \right]
 \right\},
 \tag{A.4}
\]
where the infimum is over deterministic or predictable paths in the finite
library,
\[
 \Xi_T^2=D_T^2+B^2+\sigma^2+Bb,
\]
and \(\mathcal B_{t,i}\) is the right-hand side of the simultaneous
finite-sample oracle inequality of Theorem~3.2 for the pair
\((m_i,w_{t,i})\), including discretisation and numerical optimisation
remainders.

For regular weights and polynomial-size dictionaries,
\[
 \mathcal B_{t,i}
 \lesssim
 A_0m_i^{-2\alpha}
 +\frac{b_Tm_i}{n_{\rm eff}(w_{t,i})}
 +V_{k(t)}^2A_{\gamma_{k(t)},t}(w_{t,i})^2
 +e_{m_i,T}^2+\varepsilon^{\rm opt}_{t,i}+T^{-3}.
 \tag{A.5}
\]
\end{theorem}

\paragraph{Proof.}
Allocate \(\delta/3\) to the simultaneous expert bounds, \(\delta/3\) to the
noise-envelope event, and the remainder to the bridge.  For the bridge, apply
Lemma~\ref{lem:bridge} to every path using failure probability
proportional to its prior mass \(\Pi(i_{1:T})\).  The weighted union bound gives
\(x=\mathcal C_T(i_{1:T})+\log(3/\delta)\) simultaneously over all paths.
Lemma~\ref{lem:mixable-path} controls the realised loss difference by
\(2D_T^2\mathcal C_T\).  Hence
\[
 A_T
 \le
 B_T+R_T+c_0\sqrt{x(A_T+B_T)}+c_1x.
\]
Use
\(\sqrt{A_T+B_T}\le\sqrt{A_T}+\sqrt{B_T}\) and Young's inequality with
parameter proportional to \(\zeta\).  After moving the resulting multiple of
\(A_T\) to the left,
\[
 A_T
 \le
 (1+\zeta)B_T
 +C_\zeta\Xi_T^2
 \{\mathcal C_T+\log(3/\delta)\}.
\]
Insert the simultaneous bound \(B_T\le\sum_t\mathcal B_{t,i_t}\), then take the
infimum over paths. \(\square\)

Equation (A.4) is the requested single theorem: the adaptive algorithm competes
with changing memory, kernel and network-complexity paths in the same current
risk used by the statistical paper.

\subsection{From the finite library to the continuum oracle}

Let
\[
 K_g(x)=\frac{g+1}{g}(1-x^g)\mathbf 1\{0\le x\le1\}.
\]
If the true temporal exponent is \(\gamma\), then
\[
 R_g:=\int_0^1K_g(x)^2dx=\frac{2(g+1)}{2g+1},
\]
and
\[
 M_{\gamma,g}:=\int_0^1x^\gamma K_g(x)dx
 =
 \frac{g+1}{(\gamma+1)(\gamma+g+1)}.
\]
Thus the continuum proxy associated with a kernel shape \(g\) is
\[
 F_t(m,H,g)
 =
 Am^{-2\alpha}
 +b_TR_g\frac mH
 +DV_k^2M_{\gamma_k,g}^2
 \left(\frac HT\right)^{2\gamma_k}.
 \tag{A.6}
\]
The oracle of Section~5 is the special case \(g=\gamma_k\).

Fix \(0<\gamma_{\min}\le\gamma_k\le1\) and a finite upper bound
\(\alpha\le\alpha_{\max}\).  For \(\epsilon\in(0,1)\), construct geometric
grids for \(m\) and \(H\) with ratio \(1+\epsilon\), and an additive
\(\epsilon\)-grid for \(g\in[\gamma_{\min},1]\).  Include all small integers
below \(1/\epsilon\) to remove rounding effects.

\begin{proposition}[Finite library approximates the continuum oracle]
\label{prop:continuum-grid}
Uniformly over interior triples
\(1\ll m,H\ll T\), for every
\((m,H,g)\) there is a grid point
\((\widetilde m,\widetilde H,\widetilde g)\) such that
\[
 F_t(\widetilde m,\widetilde H,\widetilde g)
 \le
 \{1+C\epsilon\}F_t(m,H,g)+o(F_t(m,H,g)),
 \tag{A.7}
\]
where \(C\) depends only on
\(\alpha_{\max},\gamma_{\min}\) and fixed envelope constants.
The library cardinality satisfies
\[
 |\mathcal I_T|
 =
 O\!\left(\epsilon^{-3}\log^2 T\right).
 \tag{A.8}
\]
For \(\epsilon_T=1/\log T\),
\[
 |\mathcal I_T|=O(\log^5T),
 \qquad
 \log|\mathcal I_T|=O(\log\log T).
\]
\end{proposition}

\paragraph{Proof.}
Choose \(\widetilde m/m\) and \(\widetilde H/H\) in
\([1,(1+\epsilon)]\), up to the explicitly included small-integer region, and
\(|\widetilde g-g|\le\epsilon\).  The three terms in (A.6) are smooth monomials
in \(m,H\).  On the compact set
\(g,\gamma\in[\gamma_{\min},1]\), both \(R_g\) and
\(M_{\gamma,g}\) are positive with uniformly bounded logarithmic derivatives.
Therefore each term changes by a factor \(1+O(\epsilon)\).  The \(o(1)\)
term is the discrete-kernel/Riemann approximation already quantified in
Section~5.2.  Counting grid points gives (A.8). \(\square\)

This proposition is the nontrivial algorithmic bridge:
\[
 \text{finite expert tracking}
 \quad\Longrightarrow\quad
 \text{adaptation to the continuum }(m,H,g)\text{ oracle}.
\]

More generally, for any continuous parameter path
\(\theta_{1:T}=(m_t,H_t,g_t)\), let
\(S_\epsilon(\theta_{1:T})\) be the number of switches of its nearest-grid
quantisation.  Combining Theorem~\ref{thm:adaptive-oracle},
Proposition~\ref{prop:continuum-grid}, and (A.2) yields
\[
 \frac1T\sum_t\mathcal E_t(\widehat f_t^{\rm A})
 \lesssim
 (1+O(\epsilon)+\zeta)
 \inf_{\theta_{1:T}}
 \frac1T\sum_tF_t(\theta_t)
 +
 \frac{C_\zeta\Xi_T^2}{T}
 \left[
 \log|\mathcal I_T|
 +S_\epsilon
 \log\!\left(\frac{e|\mathcal I_T|T}{S_\epsilon}\right)
 +\log|\mathcal Q_T|
 +\log\frac1\delta
 \right]
 +{\rm Rem}_T.
 \tag{A.9}
\]

\subsection{Recovery of the piecewise-Hölder oracle rates}

Consider the blockwise oracle construction of Section~6.4.  Inside block \(k\),
its target memory grows with the block age until it reaches the interior
\(H_{0,k}\), while its width grows according to the corresponding
memory--complexity relation.  A geometric quantisation changes either coordinate
only \(O(\epsilon^{-1}\log T)\) times per block.  Hence
\[
 S_\epsilon
 =
 O(K\epsilon^{-1}\log T).
 \tag{A.10}
\]
With \(\epsilon_T=\zeta_T=1/\log T\),
\[
 S_{\epsilon_T}=O(K\log^2T).
\]
Under the conditional Bernstein envelope,
\(\Xi_T^2=O(\log^2(T/\delta))\).  Consequently the adaptation term in
(A.9) is
\[
 {\rm AdaptCost}(K,T)
 =
 O\!\left(\frac{K\,{\rm polylog}(T/\delta)}{T}\right).
 \tag{A.11}
\]
For fixed \(K\), or more generally whenever the right-hand side is negligible
relative to the oracle statistical term, the adaptive master has exactly the
same polynomial rates as the block-aware continuum oracle, up to logarithms.

In particular, combining (A.9) with (76) gives
\[
 \frac1T\sum_{t=1}^T
 \mathcal E_t(\widehat f_t^{\rm A})
 \lesssim
 A^{1/(2\alpha+1)}
 \left(\frac{b_TK}{T}\right)^{2\alpha/(2\alpha+1)}
 +
 \frac1T\sum_{k=1}^KL_kr_{k,T}
 +
 \frac{B^2K}{T}
 +
 {\rm AdaptCost}(K,T)
 +
 {\rm Rem}_T,
 \tag{A.12}
\]
where, for \(V_k>0\),
\[
 r_{k,T}
 \asymp
 A^{\gamma_k/\Delta_k}
 b_T^{2\alpha\gamma_k/\Delta_k}
 V_k^{2\alpha/\Delta_k}
 T^{-2\alpha\gamma_k/\Delta_k},
 \qquad
 \Delta_k=\alpha+\gamma_k(2\alpha+1).
\]
Stationary blocks use the stationary boundary rate from Section~6.1.

For the Barron case \(\alpha=1/2\) and Lipschitz drift \(\gamma_k=1\), the
interior oracle term is \(T^{-2/5}\) up to logarithms, whereas the new tracking
penalty is \(T^{-1}{\rm polylog}(T)\).  The latter is asymptotically negligible.
The same conclusion holds for every fixed Sobolev smoothness in Section~5.4.

The algorithm therefore adapts, without estimating
\(\alpha,V_k,\gamma_k,\tau_k\), to the continuum memory--width oracle at the
same polynomial rate, up to logarithmic and discretisation factors.
